\documentclass[11pt,a4paper]{article}

\usepackage[T1]{fontenc}
\usepackage[utf8]{inputenc}
\usepackage[ngerman]{babel}

\usepackage{amsmath,amssymb,amsthm}
\usepackage[a4paper,margin=2.5cm]{geometry}
\usepackage{lmodern}
\usepackage{microtype}
\usepackage{hyperref}
\usepackage{enumitem}

\title{Schalenverknüpfung, \(5\)-adische Ordnung und die Rolle dichter Primzahlfünflinge}
\author{Kollaborative Ausarbeitung}
\date{\today}

\newtheorem{satz}{Satz}
\newtheorem{lemma}[satz]{Lemma}
\newtheorem{korollar}[satz]{Korollar}
\theoremstyle{definition}
\newtheorem{definition}[satz]{Definition}
\newtheorem{beispiel}[satz]{Beispiel}
\theoremstyle{remark}
\newtheorem{bemerkung}[satz]{Bemerkung}

\begin{document}
	\maketitle
	
	\begin{abstract}
		Wir untersuchen die Beziehung zwischen zwei natürlichen Schalenabbildungen, die durch
		\((2,3)\)- bzw. \((2,3,5)\)-glatte Faktoren induziert werden. Für eine natürliche Zahl \(n\)
		zeigen wir, dass die beiden schalenreduzierten Kerne eindeutig über die \(5\)-adische Ordnung
		\(v_5(n)\) miteinander verknüpft sind. Auf der Ebene der \(ABC\)-/\(E\)-Zerlegung führt dies zu
		einer exakten Intertwining-Formel der schalenreduzierten Quotienten. Weiter zeigen wir, dass auf
		reduzierter Signatur-Ebene der Übergang von der \((2,3,5)\)- zur \((2,3)\)-Schale genau als
		\(A\)-Twist erscheint. Die bereits hergeleitete Struktur dichter Primzahlfünflinge liefert dafür
		eine geometrisch-orientierte Interpretation.
	\end{abstract}
	
	\section{Einleitung}
	
	Die Zerlegung natürlicher Zahlen in glatte und nichtglatte Anteile ist in vielen arithmetischen
	Zusammenhängen ein natürlicher erster Schritt. Im hier betrachteten Rahmen tritt zusätzlich eine
	familienweise Zerlegung modulo \(12\) auf. Für Zahlen, die weder durch \(2\) noch durch \(3\)
	teilbar sind, wird eine kanonische Zerlegung in einen \(ABC\)-Anteil und einen energetischen
	\(E\)-Anteil verwendet. Dies führt zu Quotientenbildern der Form
	\[
	Q(n)=\frac{n_{ABC}}{\mathcal E(n)},
	\]
	die zunächst eher Relationen als echte Funktionen erzeugen.
	
	Eine natürliche Verfeinerung besteht darin, zunächst glatte Faktoren durch Schalenbildung
	abzutrennen. Besonders relevant sind hier die beiden Schalen
	\[
	(2,3)\qquad\text{und}\qquad (2,3,5),
	\]
	wobei der Übergang von der ersten zur zweiten Schale genau durch die \(5\)-adische Ordnung
	kontrolliert wird. Da \(5\equiv 5\pmod{12}\) zur Familie \(A\) gehört, greift diese Verfeinerung
	nicht nur äußerlich in die Faktorisierung ein, sondern direkt in die \(ABC\)-Struktur.
	
	\section{Glatte Schalen}
	
	\begin{definition}[\((2,3)\)- und \((2,3,5)\)-glatte Zahlen]
		Eine natürliche Zahl \(s\) heißt
		\begin{itemize}
			\item \((2,3)\)-glatt, wenn
			\[
			s=2^\alpha 3^\beta
			\qquad (\alpha,\beta\in\mathbb{N}_0),
			\]
			\item \((2,3,5)\)-glatt, wenn
			\[
			s=2^\alpha 3^\beta 5^\gamma
			\qquad (\alpha,\beta,\gamma\in\mathbb{N}_0).
			\]
		\end{itemize}
	\end{definition}
	
	\begin{definition}[Schalenanteile]
		Für \(n\in\mathbb{N}\) definieren wir
		\[
		s_{23}(n):=2^{v_2(n)}3^{v_3(n)},
		\qquad
		s_{235}(n):=2^{v_2(n)}3^{v_3(n)}5^{v_5(n)}.
		\]
		Die zugehörigen schalenreduzierten Kerne seien
		\[
		m_{23}(n):=\frac{n}{s_{23}(n)},
		\qquad
		m_{235}(n):=\frac{n}{s_{235}(n)}.
		\]
	\end{definition}
	
	\begin{lemma}
		Für jede natürliche Zahl \(n\) sind die Zerlegungen
		\[
		n=s_{23}(n)\,m_{23}(n)
		\qquad\text{und}\qquad
		n=s_{235}(n)\,m_{235}(n)
		\]
		eindeutig.
	\end{lemma}
	
	\begin{proof}
		Die Exponenten \(v_2(n)\), \(v_3(n)\) und \(v_5(n)\) sind durch die Primfaktorzerlegung von
		\(n\) eindeutig bestimmt. Damit sind auch \(s_{23}(n)\), \(s_{235}(n)\) und die zugehörigen Kerne
		eindeutig festgelegt.
	\end{proof}
	
	\section{Familien modulo \(12\)}
	
	\begin{definition}[Die vier Familien]
		Für Primzahlen \(p>3\) definieren wir
		\[
		E:=\{p:\ p\equiv 1 \pmod{12}\},\qquad
		A:=\{p:\ p\equiv 5 \pmod{12}\},
		\]
		\[
		B:=\{p:\ p\equiv 7 \pmod{12}\},\qquad
		C:=\{p:\ p\equiv 11 \pmod{12}\}.
		\]
	\end{definition}
	
	\begin{definition}[EABC-Gitter]
		Für \(n\in\mathbb{N}\) setzen wir
		\[
		E_n:=12n-11,\qquad
		A_n:=12n-7,\qquad
		B_n:=12n-5,\qquad
		C_n:=12n-1.
		\]
	\end{definition}
	
	\begin{definition}[Schalenreduzierte \(ABC\)-/\(E\)-Zerlegung]
		Für die Kerne \(m_{23}(n)\) und \(m_{235}(n)\) schreiben wir
		\[
		m_{23}(n)=m_{23,ABC}(n)\,\mathcal E_{23}(n),
		\]
		\[
		m_{235}(n)=m_{235,ABC}(n)\,\mathcal E_{235}(n),
		\]
		wobei der erste Faktor jeweils alle Primfaktorpotenzen aus den Familien \(A,B,C\) und der zweite
		Faktor alle Primfaktorpotenzen aus der Familie \(E\) enthält.
	\end{definition}
	
	\begin{definition}[Schalenreduzierte Quotienten]
		Wir definieren
		\[
		Q_{23}(n):=\frac{m_{23,ABC}(n)}{\mathcal E_{23}(n)},
		\qquad
		Q_{235}(n):=\frac{m_{235,ABC}(n)}{\mathcal E_{235}(n)}.
		\]
	\end{definition}
	
	\section{Algebraische Verknüpfung der beiden Schalenbilder}
	
	\begin{lemma}
		Für jede natürliche Zahl \(n\) gilt mit
		\[
		\gamma:=v_5(n)
		\]
		die Identität
		\[
		m_{23}(n)=5^\gamma\,m_{235}(n).
		\]
	\end{lemma}
	
	\begin{proof}
		Nach Definition ist
		\[
		m_{23}(n)=\frac{n}{2^{v_2(n)}3^{v_3(n)}}
		\]
		und
		\[
		m_{235}(n)=\frac{n}{2^{v_2(n)}3^{v_3(n)}5^{v_5(n)}}.
		\]
		Also folgt unmittelbar
		\[
		m_{23}(n)=5^{v_5(n)}\,m_{235}(n)=5^\gamma\,m_{235}(n).
		\]
	\end{proof}
	
	\begin{satz}[Intertwining-Formel]
		Für jede natürliche Zahl \(n\) mit \(\gamma=v_5(n)\) gilt
		\[
		m_{23,ABC}(n)=5^\gamma\,m_{235,ABC}(n),
		\]
		\[
		\mathcal E_{23}(n)=\mathcal E_{235}(n).
		\]
		Insbesondere gilt
		\[
		Q_{23}(n)=5^\gamma\,Q_{235}(n).
		\]
	\end{satz}
	
	\begin{proof}
		Da \(5\equiv 5\pmod{12}\), gehört \(5\) zur Familie \(A\), also zum \(ABC\)-Anteil und nicht zum
		\(E\)-Anteil. Beim Übergang von \(m_{23}(n)\) zu \(m_{235}(n)\) wird daher genau der Faktor
		\(5^\gamma\) aus dem \(ABC\)-Teil entfernt, während der \(E\)-Teil unverändert bleibt. Somit
		\[
		m_{23,ABC}(n)=5^\gamma\,m_{235,ABC}(n),
		\qquad
		\mathcal E_{23}(n)=\mathcal E_{235}(n).
		\]
		Die Quotientenformel folgt sofort.
	\end{proof}
	
	\section{Signaturtheoretische Interpretation}
	
	\begin{definition}[Reduzierte Signaturen]
		Sei \(\sigma_{23}(n)\in\mathbb{F}_2^3\) die reduzierte Paritätssignatur von \(m_{23,ABC}(n)\) und
		\(\sigma_{235}(n)\in\mathbb{F}_2^3\) die reduzierte Paritätssignatur von \(m_{235,ABC}(n)\).
		Wir identifizieren die \(A\)-Richtung mit
		\[
		a:=(1,0,0)\in\mathbb{F}_2^3.
		\]
	\end{definition}
	
	\begin{korollar}[Der \(A\)-Twist]
		Mit \(\gamma=v_5(n)\) gilt
		\[
		\sigma_{23}(n)=\sigma_{235}(n)+(\gamma\bmod 2)\,a.
		\]
		Insbesondere:
		\begin{enumerate}[label=\textup{(\roman*)}]
			\item Ist \(\gamma\) gerade, so gilt
			\[
			\sigma_{23}(n)=\sigma_{235}(n).
			\]
			\item Ist \(\gamma\) ungerade, so unterscheiden sich beide Schalenbilder genau um einen
			\(A\)-Schritt.
		\end{enumerate}
	\end{korollar}
	
	\begin{proof}
		Der Übergang von \(m_{235,ABC}(n)\) zu \(m_{23,ABC}(n)\) besteht in der Multiplikation mit
		\(5^\gamma\). Da \(5\) zur Familie \(A\) gehört, addiert dies auf Signaturniveau genau \(\gamma\)
		mal den Vektor \(a=(1,0,0)\). In \(\mathbb{F}_2^3\) zählt nur die Parität.
	\end{proof}
	
	\section{Dichte Primzahlfünflinge}
	
	\begin{definition}[Dichter Primzahlfünfling]
		Ein dichter Primzahlfünfling ist eine Primzahlkonstellation der Form
		\[
		(p,p+2,p+6,p+8,p+12).
		\]
	\end{definition}
	
	\begin{lemma}
		Jeder dichte Primzahlfünfling mit \(p>3\) besitzt genau eine der beiden Formen
		\[
		(A_n,B_n,C_n,E_{n+1},A_{n+1})
		\]
		oder
		\[
		(C_n,E_{n+1},A_{n+1},B_{n+1},C_{n+1})
		\]
		für ein eindeutig bestimmtes \(n\in\mathbb{N}\).
	\end{lemma}
	
	\begin{proof}
		Für \(p>3\) sind als Startreste modulo \(12\) nur \(5\) und \(11\) möglich.
		
		Falls \(p\equiv 5\pmod{12}\), so ist \(p=A_n\), also
		\[
		p+2=B_n,\quad p+6=C_n,\quad p+8=E_{n+1},\quad p+12=A_{n+1}.
		\]
		
		Falls \(p\equiv 11\pmod{12}\), so ist \(p=C_n\), also
		\[
		p+2=E_{n+1},\quad p+6=A_{n+1},\quad p+8=B_{n+1},\quad p+12=C_{n+1}.
		\]
		
		Die Fälle \(p\equiv 1\) oder \(p\equiv 7\pmod{12}\) sind unmöglich, da dann ein Folgeterm durch
		\(3\) teilbar wäre.
	\end{proof}
	
	\begin{satz}
		Die dichten Primzahlfünflinge geben der algebraischen Verknüpfung der beiden Schalenabbildungen
		eine geometrisch-orientierte Interpretation:
		
		\begin{enumerate}[label=\textup{(\roman*)}]
			\item Jeder dichte Primzahlfünfling ist eine orientierte Erweiterung einer blockübergreifenden
			Viererkonfiguration.
			\item Diese Orientierung ist genau entweder \(A\)-gerichtet oder \(C\)-gerichtet.
			\item Der Übergang von der \((2,3,5)\)-Schale zur \((2,3)\)-Schale wirkt auf reduzierter
			Signatur-Ebene als \(A\)-Twist.
			\item Damit wird die algebraische \(5\)-adische Schalenverknüpfung geometrisch durch die
			ausgezeichnete \(A\)-/\(C\)-Orientierung der Primzahlfünflinge sichtbar.
		\end{enumerate}
	\end{satz}
	
	\begin{proof}
		Die ersten beiden Aussagen folgen direkt aus dem vorangehenden Lemma. Die dritte ist genau das
		vorige Korollar. Die vierte folgt daraus, dass die zusätzliche \(5\)-adische Komponente genau in
		der \(A\)-Familie liegt und daher algebraisch dieselbe ausgezeichnete Richtung markiert, die in
		der lokalen Fünflingsgeometrie als Orientierung erscheint.
	\end{proof}
	
	\begin{korollar}
		Ohne Primzahlfünflinge ist die Formel
		\[
		Q_{23}(n)=5^{v_5(n)}\,Q_{235}(n)
		\]
		eine rein algebraische Intertwining-Relation. Durch den Einbezug der Fünflinge erhält dieselbe
		Formel eine intrinsische Orientierungsbedeutung.
	\end{korollar}
	
	\section{Beispiele}
	
	\begin{beispiel}
		Für
		\[
		n=2^3\cdot 3^2\cdot 5\cdot 7\cdot 13
		\]
		gilt
		\[
		v_5(n)=1,
		\qquad
		Q_{23}(n)=\frac{35}{13},
		\qquad
		Q_{235}(n)=\frac{7}{13}.
		\]
		Also
		\[
		Q_{23}(n)=5\,Q_{235}(n).
		\]
		Auf Signaturniveau unterscheiden sich die beiden Schalenbilder genau um einen \(A\)-Twist.
	\end{beispiel}
	
	\begin{beispiel}
		Der dichte Primzahlfünfling
		\[
		(5,7,11,13,17)
		\]
		hat die Form
		\[
		(A_1,B_1,C_1,E_2,A_2).
		\]
		Er ist also eine \(A\)-gerichtete Erweiterung. In dieser Sicht wird genau die Richtung
		hervorgehoben, die auch in der \(5\)-adischen Schalenverknüpfung ausgezeichnet ist.
	\end{beispiel}
	
	\section{Schlussbemerkung}
	
	Die \((2,3)\)- und \((2,3,5)\)-Schalenabbildungen sind algebraisch eindeutig durch die
	\(5\)-adische Ordnung miteinander verknüpft. Diese Beziehung ist nicht bloß formal:
	Da \(5\) zur Familie \(A\) gehört, erscheint der Übergang zwischen beiden Schalenbildern auf
	Signaturniveau als \(A\)-Twist. Die dichten Primzahlfünflinge liefern dafür die natürliche lokale
	Geometrie, denn sie sind genau die ersten Konstellationen, in denen eine ausgezeichnete
	\(A\)- oder \(C\)-Orientierung sichtbar wird.
	
	Damit entsteht eine kohärente Brücke zwischen
	\begin{enumerate}[label=\arabic*.]
		\item glatten Schalen,
		\item \(ABC\)-/\(E\)-Moden,
		\item reduzierten Signaturen,
		\item und der lokalen Geometrie dichter Primzahlkonstellationen.
	\end{enumerate}
	
\end{document}